\begin{problemsummary}
{A composition error names a query you did not write, through a service you did
not touch, and stops at a field that looks correct. Or the graph composes at
exit zero and a field answers HTTP 500, or answers a different field's value,
or comes back null with no \code{errors} key at all.}
{The check asks whether a field can be reached at all, not whether each route
to it works, so it reports one route out of however many fail and picks the one
it happened to be walking. The route it picks moves with the rest of the graph:
which other subgraphs are in \code{graph.yaml}, the order their lines are
written in, and the order root fields are declared inside one document all
change the message without changing the defect. And the check can be satisfied
without finding a route at all: a root field returning an interface, declared
\code{@shareable} in two subgraphs, is accepted when either subgraph answers
for its own types, which records the entity as visited and waves every later
route to it through.}
{Read the second reason line first, because it names the subgraph to open, and
treat the printed path and the entity ancestor's subgraph as working rather
than as diagnosis. Put \code{resolvable: false} on a key only where the type
carries nothing but that key: a subgraph contributing one field to a type will
be sent representations of it and has to be enterable, whether or not it owns a
row. Then stop treating a pass as proof. Compose each contributing subgraph
against the one service it hangs its fields on, and send one request per
contributed field through the router, which is the only check in this book that
catches all three of the failures above.}
\end{problemsummary}
